\documentclass[10pt, aspectratio=169]{beamer}
% Preámbulo modular para presentaciones Beamer (Modelación Estadística)
% Diseño y paleta institucional del Tecnológico de Monterrey

\documentclass[aspectratio=169,xcolor={dvipsnames,table}]{beamer}

\usepackage[utf8]{inputenc}
\usepackage[spanish,mexico]{babel}
\usepackage{amsmath,amssymb,amsfonts,mathtools}
\usepackage{graphicx}
\usepackage{listings}
\usepackage{booktabs}
\usepackage{tikz}
\usetikzlibrary{matrix,arrows,decorations.pathmorphing,shapes.geometric,positioning}

% Paleta Institucional Tecnológico de Monterrey
\definecolor{TecAzul}{HTML}{0039A6}       % Azul TEC: color primario institucional
\definecolor{TecAzulOscuro}{HTML}{1A2E51} % Azul marino oscuro
\definecolor{TecRojo}{HTML}{EC2661}       % Rojo de acento (paleta miTec)
\definecolor{TecGris}{HTML}{646464}       % Gris neutro para comentarios y subtítulos
\definecolor{TecFondoBloque}{HTML}{F4F6F9}% Fondo suave para bloques

% Configuración de Tema Beamer
\usetheme{Boadilla}
\usecolortheme[named=TecAzulOscuro]{structure}
\setbeamercolor{alerted text}{fg=TecRojo}
\setbeamercolor{palette primary}{bg=TecAzulOscuro,fg=white}
\setbeamercolor{palette secondary}{bg=TecAzul,fg=white}
\setbeamercolor{palette tertiary}{bg=TecAzulOscuro!80!black,fg=white}
\setbeamercolor{title}{fg=TecAzulOscuro,bg=TecFondoBloque}
\setbeamercolor{frametitle}{fg=TecAzulOscuro,bg=TecFondoBloque}

% Bloques personalizados elegantes
\setbeamercolor{block title}{bg=TecAzulOscuro,fg=white}
\setbeamercolor{block body}{bg=TecFondoBloque,fg=black}
\setbeamercolor{block title alerted}{bg=TecRojo,fg=white}
\setbeamercolor{block body alerted}{bg=TecRojo!8,fg=black}
\setbeamercolor{block title example}{bg=TecAzul,fg=white}
\setbeamercolor{block body example}{bg=TecAzul!8,fg=black}

% Redondear bordes y añadir sombra sutil
\setbeamertemplate{blocks}[rounded][shadow=true]
\setbeamertemplate{navigation symbols}{}

% Pie de página limpio con número de diapositiva
\setbeamertemplate{footline}{
  \leavevmode%
  \hbox{%
  \begin{beamercolorbox}[wd=.7\paperwidth,ht=2.25ex,dp=1ex,left]{author in head/foot}%
    \hspace*{2ex}\insertshortauthor~~(\insertshortinstitute) \hfill \insertshorttitle
  \end{beamercolorbox}%
  \begin{beamercolorbox}[wd=.3\paperwidth,ht=2.25ex,dp=1ex,right]{date in head/foot}%
    \insertshortdate{}\hspace*{2em}
    \insertframenumber{} / \inserttotalframenumber\hspace*{2ex}
  \end{beamercolorbox}}%
  \vskip0pt%
}

% Configuración de Listings de Python para visualización pedagógica de código
\lstset{
    language=Python,
    basicstyle=\ttfamily\footnotesize,
    keywordstyle=\color{TecAzulOscuro}\bfseries,
    stringstyle=\color{TecRojo},
    commentstyle=\color{TecGris}\itshape,
    showstringspaces=false,
    breaklines=true,
    frame=lines,
    rulecolor=\color{TecAzulOscuro!30},
    backgroundcolor=\color{TecFondoBloque},
    aboveskip=0.5em,
    belowskip=0.5em,
    literate={á}{{\'a}}1 {é}{{\'e}}1 {í}{{\'i}}1 {ó}{{\'o}}1 {ú}{{\'u}}1
             {Á}{{\'A}}1 {É}{{\'E}}1 {Í}{{\'I}}1 {Ó}{{\'O}}1 {Ú}{{\'U}}1
             {ñ}{{\~n}}1 {Ñ}{{\~N}}1 {¿}{{?`}}1 {¡}{{!`}}1
}

% Metadatos institucionales por defecto
\title[Modelación Estadística]{Modelación Estadística para Ciencia de Datos e Ingeniería}
\author[J. Castillo Colmenares]{Juliho David Castillo Colmenares}
\institute[Tec de Monterrey]{Tecnológico de Monterrey \\ Escuela de Ingeniería y Ciencias}
\date{\today}

% Comandos y macros estadísticas compartidas para las presentaciones Beamer (Metropolis)

% Conjuntos numéricos básicos
\newcommand{\R}{\mathbb{R}}
\newcommand{\N}{\mathbb{N}}
\newcommand{\Q}{\mathbb{Q}}
\newcommand{\Z}{\mathbb{Z}}
\newcommand{\set}[1]{\left\{ #1 \right\}}
\newcommand{\abs}[1]{\left|#1\right|}
\newcommand{\norm}[1]{\left\| #1 \right\|}

% Comandos estadísticos y probabilísticos del curso
\newcommand{\Var}{\operatorname{Var}}
\newcommand{\cov}{\operatorname{Cov}}
\newcommand{\corr}{\rho}
\newcommand{\comb}[2]{\begin{pmatrix} #1 \\ #2 \end{pmatrix}}
\newcommand{\s}{\sigma}
\newcommand{\particion}{\mathfrak{P}}
\newcommand{\card}[1]{\#\left( #1 \right)}
\newcommand{\seq}[3]{\set{#1}_{#2}^{#3}}
\newcommand{\E}{\mathbb{E}}
\newcommand{\Prob}{\mathbb{P}}

% Entornos pedagógicos y cajas en español académico natural
\newenvironment{discusion}{%
  \begin{alertblock}{Pregunta de Discusión}%
}{%
  \end{alertblock}%
}

\newenvironment{reflexion}{%
  \begin{alertblock}{Reflexión para el Aula}%
}{%
  \end{alertblock}%
}

\newenvironment{paradoja}{%
  \begin{alertblock}{Paradoja de Exploración}%
}{%
  \end{alertblock}%
}

\newenvironment{motivacion}{%
  \begin{exampleblock}{Motivación: Ciencia de Datos en la Práctica}%
}{%
  \end{exampleblock}%
}

\newenvironment{retoclase}{%
  \begin{block}{Reto Rápido en Equipo}%
}{%
  \end{block}%
}


\title[Coeficientes Óptimos]{Sección 09.05: Coeficientes Óptimos y Pruebas de Significancia}
\subtitle{Prueba $t$ para la Pendiente, la Identidad $F=t^2$, y el Error Estándar Residual}
\author[J. Castillo Colmenares]{Juliho Castillo Colmenares}
\institute{Tecnológico de Monterrey}
\date{\vspace{-1.2cm}}

\begin{document}

% Slide 1: Portada
\begin{frame}[plain]
  \vspace{-0.8cm}
  \titlepage
\end{frame}

% Slide 2: Hoja de Ruta
\begin{frame}{Hoja de Ruta --- Capítulo 09: Regresiones Lineales y Múltiples}
  \scriptsize
  ¿Dónde nos encontramos en el desarrollo modular del capítulo?
  \begin{itemize}\itemsep=0.02em
    \item \textbf{09.01} Correlación como Premisa de la Regresión
    \item \textbf{09.02} Introducción a la Regresión Lineal
    \item \textbf{09.03} Matemáticas de la Regresión: Derivación MCO y $R^2$
    \item \textbf{09.04} Regresión Lineal sobre Datos Simulados
    \item \textbf{09.05} Coeficientes Óptimos, Pruebas $t$/$F$ y RSE \emph{(Hoy)}
    \item \textbf{09.06} Implementación en Python con \texttt{statsmodels}
    \item \textbf{09.07} Regresión Múltiple, Ecuación Normal y Ridge/Lasso
    \item \textbf{09.08} Validación de Modelos y $k$-fold Cross-Validation
    \item \textbf{09.09} Diagnóstico de Residuos y Multicolinealidad (VIF)
    \item \textbf{09.10} Regresión con \texttt{scikit-learn}
    \item \textbf{09.11} Variables Categóricas y Variables Muda
    \item \textbf{09.12} Transformaciones No Lineales y Regresión Polinomial
  \end{itemize}
  \pause
  \vspace{-0.05cm}
  \begin{block}{Objetivo de la Sesión}
    Someter el coeficiente de pendiente estimado a una prueba de hipótesis formal, verificar la identidad algebraica $F=t^2$ que conecta la significancia individual con la global en regresión simple, y calcular el Error Estándar Residual (RSE) como medida absoluta --- no relativa --- de la calidad del ajuste.
  \end{block}
\end{frame}

% Slide 3: Motivación
\begin{frame}{¿Son Confiables los Coeficientes Estimados?}
  \footnotesize
  \begin{motivacion}
    Al calcular los valores de $\hat\alpha$ y $\hat\beta$, obtenemos \textbf{estimados}, no valores exactos: cada uno lleva asociado un valor-$p$. Necesitamos demostrar su significancia estadística mediante una prueba de hipótesis antes de confiar en el modelo.
  \end{motivacion}

  \vspace{0.15cm}
  \pause
  La pregunta central: ¿el valor $\hat\beta$ que obtuvimos es diferente de cero, o podría ser simplemente ruido muestral alrededor de $\beta=0$ (ausencia real de relación lineal)?
\end{frame}

% Slide 4: Prueba de hipótesis para beta
\begin{frame}{Prueba de Hipótesis para la Pendiente}
  \footnotesize
  \begin{block}{Planteamiento}
    $$H_0:\beta=0 \quad (\text{no hay relación lineal}) \qquad \text{vs.} \qquad H_a:\beta\neq0$$
  \end{block}
  \pause

  \vspace{0.1cm}
  \begin{alertblock}{Estadístico de Prueba}
    $$t=\frac{\hat\beta}{SE(\hat\beta)}, \qquad SE(\hat\beta)=\sqrt{\frac{\text{CME}}{S_{XX}}} \ \sim\ t_{n-2} \quad\text{bajo } H_0$$
  \end{alertblock}
\end{frame}

% Slide 5: La identidad F=t²
\begin{frame}{La Identidad $F=t^2$ en Regresión Simple}
  \footnotesize
  \begin{block}{Prueba $F$ Global}
    $$F=\frac{\text{SCR}/1}{\text{SCE}/(n-2)} \sim F_{1,n-2}$$
  \end{block}
  \pause

  \vspace{0.1cm}
  \begin{alertblock}{Conexión Exacta con la Prueba $t$}
    En regresión simple ($p=1$ predictor), la prueba $F$ global de significancia del modelo y la prueba $t$ individual de la pendiente son \textbf{matemáticamente equivalentes}: $F=t^2$. Con múltiples predictores (Sección 09.07), esta equivalencia desaparece: el estadístico $F$ evalúa la significancia \emph{conjunta}, mientras que cada $t_j$ evalúa una significancia \emph{individual}.
  \end{alertblock}
\end{frame}

% Slide 6: Error Estándar Residual
\begin{frame}{El Error Estándar Residual (RSE)}
  \footnotesize
  \begin{block}{Definición}
    $$RSE=\sqrt{\frac{\text{SCE}}{n-2}}$$
  \end{block}
  \pause

  \vspace{0.1cm}
  \begin{alertblock}{Interpretación}
    RSE estima la desviación estándar del término de error: es el error de predicción típico \emph{en las unidades originales de $Y$}, incluso si los coeficientes fueran conocidos con exactitud absoluta. Se interpreta comparándolo contra $\bar Y$: $RSE/\bar Y$ da un error relativo porcentual.
  \end{alertblock}
\end{frame}

% Slide 7: Lab Python (Bloque 1)
\begin{frame}[fragile]{Laboratorio en Python: Coeficientes Óptimos y Prueba $t$}
  \scriptsize
  Ajuste MCO y prueba de significancia de la pendiente (\texttt{09.05\_optimal\_coefficients\_and\_tests.py}):
  {\lstset{basicstyle=\fontsize{3.4pt}{4.1pt}\selectfont\ttfamily,aboveskip=0.05em,belowskip=0.05em}
  \lstinputlisting[firstline=18, lastline=35]{../../code/09_regresiones/09.05_optimal_coefficients_and_tests.py}}
\end{frame}

% Slide 8: Lab Python (Bloque 2)
\begin{frame}[fragile]{Laboratorio en Python: Verificación de $F=t^2$}
  \scriptsize
  Confirmación numérica de la identidad algebraica exacta entre las pruebas $t$ y $F$:
  {\lstset{basicstyle=\fontsize{3.6pt}{4.3pt}\selectfont\ttfamily,aboveskip=0.05em,belowskip=0.05em}
  \lstinputlisting[firstline=41, lastline=51]{../../code/09_regresiones/09.05_optimal_coefficients_and_tests.py}}
\end{frame}

% Slide 9: Lab Python (Bloque 3)
\begin{frame}[fragile]{Laboratorio en Python: Error Estándar Residual}
  \scriptsize
  Cálculo del RSE y su interpretación como error relativo:
  {\lstset{basicstyle=\fontsize{3.8pt}{4.5pt}\selectfont\ttfamily,aboveskip=0.05em,belowskip=0.05em}
  \lstinputlisting[firstline=54, lastline=63]{../../code/09_regresiones/09.05_optimal_coefficients_and_tests.py}}
\end{frame}

% Slide 10: Salida Terminal
\begin{frame}[fragile]{Laboratorio en Python: Salida en Terminal}
  \scriptsize
  Salida estándar generada al ejecutar el script del laboratorio en Python:
  \vspace{0.02cm}
  \begin{block}{Salida Estándar en Consola}
    \fontsize{3.6pt}{4.3pt}\selectfont
    \begin{verbatim}
=== Block 1: Optimal Coefficients and t-Test ===
alpha_hat=1.8577, beta_hat=0.3468
SE(beta_hat)=0.0202, t=17.1624 (df=98), p-value=2.73e-31

=== Block 2: F-Test (F = t^2) ===
F statistic = 294.5482 (df=1,98), p-value=2.73e-31
t^2 = 294.5482, |F - t^2| = 0.0000000000

=== Block 3: Residual Standard Error ===
RSE = 0.4453
mean(Y) = 2.4203
Relative error = RSE/mean(Y) = 0.1840 (18.4%)
    \end{verbatim}
  \end{block}
\end{frame}

% Slide 13: Cierre
\begin{frame}{Cierre de Sección: Coeficientes Óptimos y Pruebas de Significancia}
  \begin{reflexion}
    Un coeficiente estimado sin una prueba de significancia es solo un número: hemos formalizado cuándo $\hat\beta$ representa evidencia real de una relación lineal, descubierto la elegante identidad $F=t^2$ propia de la regresión simple, y establecido el RSE como la medida absoluta de referencia para juzgar la calidad predictiva de cualquier modelo.
  \end{reflexion}

  \vspace{0.2cm}
  \pause
  \begin{block}{Lecciones Clave}
    \begin{itemize}\itemsep=0.08cm
      \item \textbf{Prueba $t$:} $t=\hat\beta/SE(\hat\beta)\sim t_{n-2}$ bajo $H_0:\beta=0$.
      \item \textbf{$F=t^2$:} equivalencia exacta exclusiva de la regresión de un solo predictor.
      \item \textbf{RSE:} error absoluto de predicción en las unidades de $Y$; $RSE/\bar Y$ da el error relativo.
    \end{itemize}
  \end{block}
\end{frame}

% Slide 14: Perspectiva Modular
\begin{frame}{Perspectiva Modular --- El Próximo Reto}
  \scriptsize
  \begin{retoclase}
    Hemos calculado todo manualmente con fórmulas cerradas. La Sección 09.06 muestra cómo obtener exactamente los mismos resultados --- coeficientes, valores-$p$, $R^2$, RSE --- con una sola línea de código profesional usando la librería \texttt{statsmodels}, la herramienta estándar de la industria para regresión estadística en Python.
  \end{retoclase}

  \vspace{0.08cm}
  \pause
  \begin{alertblock}{Preparación Recomendada}
    \begin{itemize}\itemsep=0.06cm
      \item Instala \texttt{statsmodels} si aún no lo tienes (\texttt{pip install statsmodels}).
      \item Reflexiona sobre las ventajas y riesgos de usar una librería "caja negra" vs. implementar las fórmulas manualmente como en esta sección.
    \end{itemize}
  \end{alertblock}
\end{frame}

\end{document}
